% ======================================================================
% CAPÍTULO 4 — METODOLOGIA
% ======================================================================
\chapter{METODOLOGIA}

Este capítulo descreve a metodologia adotada no desenvolvimento do sistema,
abrangendo a natureza da pesquisa, a estratégia de desenvolvimento, os materiais
utilizados, os procedimentos executados, o ambiente de validação e os critérios
de aceitação. A descrição segue o princípio de distinguir o que foi planejado do
que foi efetivamente executado.

\section{Natureza e abordagem da pesquisa}

O trabalho caracteriza-se como uma pesquisa aplicada, de natureza experimental,
com abordagem predominantemente quantitativa. O objetivo é gerar um artefato
tecnológico --- um sistema de automação e supervisão --- e avaliá-lo quanto ao
atendimento de requisitos funcionais e de segurança. A estruturação do texto
segue as normas de documentação de trabalhos acadêmicos \cite{abnt14724,
abnt6023, abnt10520}.

A estratégia metodológica combina duas práticas de engenharia de software:
(i) a \emph{especificação dirigida}, em que os requisitos e os contratos de
interface são consolidados antes da implementação; e (ii) o \emph{desenvolvimento
orientado a testes} (TDD), em que os critérios de aceitação são traduzidos em
testes automatizados antes da codificação de cada funcionalidade
\cite{beck2002}. Essa combinação visa garantir a rastreabilidade entre
requisitos, implementação e validação.

\section{Etapas do desenvolvimento}

O desenvolvimento foi organizado nas seguintes etapas:

\begin{enumerate}[label=\arabic*)]
  \item \textbf{Levantamento e análise de requisitos:} consolidação da
  documentação técnica do processo (especificação do sistema e requisitos de
  automação), identificação das figuras de referência do processo, da matriz de
  acionamento dos atuadores e dos parâmetros de controle;
  \item \textbf{Definição da arquitetura:} escolha da topologia mestre-escravo,
  definição dos módulos de software e firmware e dos contratos de interface
  (protocolo serial e API);
  \item \textbf{Desenvolvimento incremental do firmware:} implementação e teste
  da camada \daq --- acionamento de atuadores, leitura de termopares, protocolo
  JSON e watchdog;
  \item \textbf{Desenvolvimento incremental do backend:} implementação da FSM,
  dos controladores PID e da rampa, da persistência de parâmetros, do loop de
  controle a 4~Hz e da API REST/WebSocket;
  \item \textbf{Desenvolvimento incremental da IHM Web:} implementação dos
  componentes de monitoramento, do Diagrama de Tempos, dos gráficos de tendência
  e dos painéis de controle e configuração;
  \item \textbf{Validação:} execução de testes unitários, de integração,
  funcionais e de segurança, com correção iterativa.
\end{enumerate}

Cada etapa foi conduzida de forma incremental, com validação contínua por meio de
testes automatizados, e os contratos de dados (JSON de escrita/leitura,
telemetria) foram definidos antecipadamente e utilizados como especificação de
interface entre os módulos.

\section{Materiais e ambiente}

\subsection{Hardware de referência}

O processo físico de referência é composto por: frasco de reação com agitação
por gás hélio; bomba peristáltica para injeção de TEBS; v\'alvulas solen\'oides
SV1 a SV5; dissecante de n\'afion; Tubo U com filamento resistivo (Forno 1) e
copo de nitrog\^enio l\'iquido com pist\~ao criog\^enico; Forno 2 (atomizador);
termopares tipo K; e detector Lumex. O acionamento das v\'alvulas e da bomba \'e
intermediado por m\'odulos de rel\'e, e o acionamento dos fornos, por rel\'es de
estado s\'olido \cite{arduino, lumex}.

A plataforma de controle \'e composta por um Raspberry Pi (execu\c{c}\~ao do
backend e da IHM) e por um Arduino Uno (execu\c{c}\~ao do \daq de campo). Os
termopares s\~ao lidos por conversores do tipo MAX6675/MAX31855 via SPI
\cite{max6675}. A identifica\c{c}\~ao exata do circuito integrado empregado no
hardware \'e um item a confirmar com a documenta\c{c}\~ao f\'isica do
equipamento.

\subsection{Software e ferramentas}

A Tabela~\ref{tab:materiais} resume o ambiente de software utilizado.

\begin{table}[htbp]
\centering
\caption{Ambiente de software e ferramentas empregados no desenvolvimento.}
\label{tab:materiais}
\begin{tabular}{lll}
\toprule
\textbf{Componente} & \textbf{Tecnologia} & \textbf{Função} \\
\midrule
Backend & Python 3.11+; FastAPI & Lógica de supervisão, FSM, PID, API/WS \\
Firmware & C++; PlatformIO; ArduinoJson & Firmware do DAQ (Arduino Uno) \\
IHM Web & React; TypeScript; Vite & Interface de supervisão no navegador \\
Testes & pytest; PlatformIO test; Vitest & Testes automatizados das camadas \\
Simulação & socat + simulador do DAQ & Enlace serial virtual sem hardware \\
\bottomrule
\end{tabular}
\par\vspace{2pt}{\footnotesize Fonte: elaborada pelo autor com base na documenta\c{c}\~ao t\'ecnica do projeto.}
\end{table}

\section{Ambiente de valida\c{c}\~ao}

A valida\c{c}\~ao foi realizada em ambiente de desenvolvimento, sem o hardware
f\'isico do processo, por tr\^es estrat\'egias complementares:

\begin{itemize}
  \item \textbf{Testes unit\'arios:} valida\c{c}\~ao isolada de cada m\'odulo
  (FSM, PID, rampa, persist\^encia, parser JSON do firmware, watchdog);
  \item \textbf{Testes de integra\c{c}\~ao e funcionais:} valida\c{c}\~ao da
  comunica\c{c}\~ao serial a 4~Hz e do ciclo completo $T_0 \to T_3$, verificando
  a matriz de atuadores em cada fase. Para tanto, um par de portas seriais
  virtuais \'e criado com \texttt{socat}, conectando o backend a um simulador do
  \daq que emula o protocolo JSON;
  \item \textbf{Testes de seguran\c{c}a:} simula\c{c}\~ao de perda de
  comunica\c{c}\~ao para confirmar a atua\c{c}\~ao do watchdog e o retorno ao
  Safe State.
\end{itemize}

O uso de um simulador de hardware permitiu exercitar o ciclo de controle completo
e validar o comportamento temporal e a seguran\c{c}a sem depender da bancada
f\'isica, o que \'e adequado \`a fase de desenvolvimento de software. A
valida\c{c}\~ao em bancada com o processo real, incluindo a calibra\c{c}\~ao da
curva de aquecimento e a verifica\c{c}\~ao experimental da rampa, permanece como
etapa de continuidade.

\section{Vari\'aveis, par\^ametros e crit\'erios de aceita\c{c}\~ao}

As principais vari\'aveis de processo s\~ao as temperaturas do Tubo U ($T_1$) e
do Forno 2 ($T_2$); as vari\'aveis manipuladas s\~ao as sa\'idas PWM ($u$ e
$f_2$); e as sa\'idas discretas s\~ao as v\'alvulas SV1 a SV5 e a bomba. Os
par\^ametros configur\'aveis do m\'etodo s\~ao os ganhos PID de cada forno, os
tempos $t_1$, $t_2$ e $t_3$, o tempo de rampa, a temperatura do nitrog\^enio, a
temperatura alvo do Tubo U e o setpoint do Forno 2.

A Tabela~\ref{tab:criterios} consolida os crit\'erios de aceita\c{c}\~ao
definidos para a valida\c{c}\~ao do sistema, conforme a documenta\c{c}\~ao de
requisitos.

\begin{table}[htbp]
\centering
\caption{Crit\'erios de aceita\c{c}\~ao e m\'etricas de valida\c{c}\~ao do sistema.}
\label{tab:criterios}
\begin{tabular}{p{4.6cm}p{4.6cm}p{5.4cm}}
\toprule
\textbf{M\'etrica} & \textbf{Alvo} & \textbf{Forma de medi\c{c}\~ao} \\
\midrule
Linearidade da rampa (Tubo U) & Desvio $\le \pm 2\,\celsius$ do SP & Curva VP $\times$ SP em bancada \\
Estabilidade do Forno 2 & $\pm 5\,\celsius$ em torno de 700\,\celsius & Curva VP $\times$ SP em bancada \\
Disponibilidade da comunica\c{c}\~ao & 100\% em 4~h de estresse a 4~Hz & Teste de estresse integrado \\
Persist\^encia de par\^ametros & 100\% \'integra ap\'os rein\'icio & Teste funcional \\
Seguran\c{c}a (watchdog) & Fornos desligados em $\le 1$~s & Teste de perda de serial \\
\bottomrule
\end{tabular}
\par\vspace{2pt}{\footnotesize Fonte: elaborada pelo autor com base na documenta\c{c}\~ao de requisitos do sistema.}
\end{table}

\section{Planejado versus executado}

No planejamento, a valida\c{c}\~ao seria conclu\'ida com a verifica\c{c}\~ao em
bancada das curvas reais de temperatura e da calibra\c{c}\~ao da curva de
aquecimento do Tubo U. Na execu\c{c}\~ao desta fase do trabalho, a valida\c{c}\~ao
foi realizada integralmente em ambiente simulado (testes automatizados), tendo
em vista a indisponibilidade da bancada f\'isica no per\'iodo de desenvolvimento.
Essa limita\c{c}\~ao \'e explicitada e registrada como etapa de continuidade,
n\~ao comprometendo a valida\c{c}\~ao da l\'ogica de controle e dos mecanismos de
seguran\c{c}a, que independem do hardware para sua correta execu\c{c}\~ao.
